In a small invertible sequent calculus, a transformer tactic policy trained to imitate optimal tactics on formulas with \rhval{const/train_lo}--\rhval{const/train_hi} connectives guides best-first search better than breadth-first search on larger formulas (H1 supported) but not better than a tuned hand heuristic (H2 refuted), and its deficit does not shrink in a monotone fashion with size (H3 refuted). Searching on cumulative $-\log p$ is what makes the transformer lose most: the same model decoded greedily or ranked reaches the heuristic's level at 41--70 connectives (H4 refuted: search with this cost is worse than greedy). Training on fewer sequents (1000 instead of \rhval{const/n_train_seq}) improved extrapolation, an unexplained exploratory result. A feature MLP with a rank cost exceeded the heuristic in a post-hoc comparison, which we offer as a hypothesis for a larger study. Natural next steps are a pre-registered test of rank versus likelihood costs with more seeds, policies that predict remaining proof size, non-invertible calculi where wrong tactics have a cost, and a measurement of proof-size optimality beyond the solved-within-budget metric.
